\chapter{Vérifications et dépannage}
\section{Une commande paraît indisponible}
Vérifier qu’une photo est ouverte, attendre la fin du traitement et contrôler
qu’il s’agit d’un RAW pour une commande de développement appareil. Un profil
absent de la liste n’est pas nécessairement une erreur de la liste : sa
disponibilité dépend des fichiers installés et de la compatibilité.

\section{La photo semble différente d’un autre logiciel}
Comparer d’abord le même fichier, le même cadrage et le même état neutre.
Vérifier le profil appareil, la balance des blancs, la calibration, les corrections
optiques, le bruit et la netteté. Une même température affichée n’établit pas
à elle seule une égalité de développement. Le fond de visualisation et le
chemin de gestion couleur peuvent également influencer la comparaison.

\section{Des réglages réapparaissent à l’ouverture}
Rechercher le compagnon \chemin{.ipr} conservé avec la photo. Sa restauration
est volontaire. Utiliser la réinitialisation adaptée à l’intention, plutôt
que supprimer immédiatement des fichiers. Conserver une copie avant toute
intervention manuelle sur un compagnon ou une configuration.

\section{Une sauvegarde ou un export échoue}
Lire le message au bas de la fenêtre. Vérifier l’espace disque, les droits
du dossier, sa disponibilité et les éventuels verrouillages. Un message
d’échec de compagnon n’est pas un export réussi. Le message de réussite
d’export indique le chemin réel, y compris le numéro ajouté si nécessaire.

\section{Une fonction n’est pas décrite}
Les panneaux recadrage, masques et IA sont réservés dans cette édition.
Les badges complets du ruban, l’état exhaustif des groupes modifiés et les
états détaillés du moteur restent des évolutions, pas des commandes cachées.
La validation visuelle par l’utilisateur et les tests sur d’autres boîtiers
restent nécessaires malgré les tests automatisés de l’IHM.
